\section*{Erd\H{o}s Problem 900: long paths in a sparse random graph}

\paragraph{Statement and result.}
For every fixed \(c>1/2\), the uniform random graph
\(G(n,\lfloor cn\rfloor)\) has with probability tending to one
a path of at least \(f(c)n\) edges, for a positive function with
\[
 f(c)\longrightarrow0\ (c\downarrow1/2),\qquad
 f(c)\longrightarrow1\ (c\to\infty).
\]
We reconstruct a depth-first-search proof for the first
claim and give a separate elementary estimate for large \(c\).
Standard binomial concentration is the only probabilistic
estimate used beyond union bounds.

\paragraph{Two search invariants.}
Depth-first search partitions the vertices into a finished
set \(S\), a stack \(U\), and an undiscovered set \(T\).
The stack is always a simple path. Every pair in \(S\times T\)
has already been queried and found to be a nonedge.
When the stack is empty, an undiscovered vertex starts a
new search. Every positive answer to an edge query discovers
a new vertex, so the number of discoveries \(|S|+|U|\)
is at least the number of positive answers.
In \(G(n,p)\), the answers to successive previously unqueried
edges are independent Bernoulli variables with parameter \(p\),
even though the search chooses its queries adaptively.

\paragraph{A quantitative supercritical bound.}
Fix \(0<\epsilon\leq1/8\), set \(p=(1+\epsilon)/n\),
\(u=\epsilon^2n/16\), and \(t_0=\lfloor\epsilon n^2/2\rfloor\).
Suppose, towards a contradiction, that the search stack
always has size at most \(u\).
By the first time \(|S|=\lceil\epsilon n\rceil\), the number
of queries is at least
\[
 |S||T|\geq(\epsilon n+O(1))
               (n-\epsilon n-u-O(1))>t_0.
\]
Thus the search makes \(t_0\) queries, and at that time
\(|S|\leq\epsilon n+O(1)\).
The sum of their Bernoulli answers has mean
\((\epsilon/2+\epsilon^2/2)n+O(1)\). A Chernoff bound gives,
with probability tending to one, at least
\((\epsilon/2+3\epsilon^2/8)n\) positive answers.
Hence at that time
\[
 |S|\geq s_0n,\qquad s_0=\epsilon/2+5\epsilon^2/16.
\]
For \(s\leq\epsilon n+O(1)\) the function \(s(n-s-u)\)
is increasing, so the already queried pairs in \(S\times T\)
number at least
\[
 s_0(1-s_0-\epsilon^2/16)n^2
 =\left(\frac\epsilon2+\frac{\epsilon^2}{16}
       -\frac{11\epsilon^3}{32}
       -\frac{15\epsilon^4}{128}\right)n^2
 >t_0.
\]
The strict inequality holds for every \(0<\epsilon\leq1/8\).
This contradicts the definition of \(t_0\).
Therefore a stack of more than \(u\) vertices occurs,
and for sufficiently large \(n\) it gives a path with
at least \(\epsilon^2n/32\) edges.

\paragraph{Transfer to a fixed number of edges.}
For given \(c>1/2\), choose
\[
 \epsilon(c)=\min((2c-1)/2,1/8).
\]
Then \((1+\epsilon(c))/2<c\).
Choose a uniform random ordering of all possible edges
and an independent binomial number \(K\) of edges with
parameters \(\binom n2,(1+\epsilon(c))/n\).
The first \(K\) edges have distribution \(G(n,p)\);
the first \(\lfloor cn\rfloor\) have the uniform fixed-size
distribution. Concentration gives
\(\Pr(K\leq\lfloor cn\rfloor)\to1\), so the path just proved
persists in the latter graph.

\paragraph{A path occupying nearly all vertices when \(c\) is large.}
Put \(\alpha=c^{-1/3}\) and \(r=\lceil\alpha n\rceil\).
For fixed disjoint \(r\)-sets \(X,Y\), the probability that
the uniform graph has no edge between them is at most
\[
 \left(1-\frac{r^2}{\binom n2}\right)^{\lfloor cn\rfloor}
 \leq \exp(-(2c\alpha^2+o(1))n).
\]
There are at most \(3^n\) ordered disjoint pairs of subsets.
For fixed \(c\geq64\), the union bound therefore shows that
every such pair has a connecting edge, with high probability.
In depth-first search, consider the instant \(|S|=r\).
There can then be fewer than \(r\) vertices in \(T\), since
there are no edges from \(S\) to \(T\).
The stack has more than \(n-2r\) vertices. For large \(n\)
it gives a path with at least \((1-3c^{-1/3})n\) edges.

\paragraph{An explicit admissible function and attribution.}
Take
\[
 f(c)=
 \begin{cases}
  \epsilon(c)^2/32,&1/2<c<64,\\
  1-3c^{-1/3},&c\geq64.
 \end{cases}
\]
It has both required limits and is positive everywhere.
The original resolution is due to Ajtai, Koml\'os and
Szemer\'edi. The search method above follows Krivelevich
and Sudakov, \emph{The phase transition in random graphs---a
simple proof}, \url{https://arxiv.org/abs/1201.6529},
with conservative constants and the model-transfer argument
written out. See \url{https://www.erdosproblems.com/900},
checked 24 September 2026. This is a complete reconstruction,
not a novel result or a local formal verification.

\paragraph{Completion estimate.}
\textbf{COMPLETION: 100\%}
